% Generated by prepare_report_f3.py from the run's analysis tables
\begin{itemize}
  \item \textbf{Income decides the resource.} The van share is 0.9\,\% in the lowest and 31.7\,\% in the highest income band; van users have a mean income of 0.74 and FCS users 0.48.
  \item \textbf{Low-income customers} (0.0-0.2) are served by a van in 0.0\,\% of the peak and 1.3\,\% of the off-peak cases; they stay at the FCS in both periods.
  \item \textbf{Peak vs off-peak.} In the top income band the van share is 26.3\,\% in the peak against 34.1\,\% off-peak (lower in the peak); in the peak the vans are busiest, so fewer high-income customers can be fitted in.
  \item \textbf{Planning horizon.} In the top income band the mean van share is 23.6\,\% in the short-horizon (type 1) and 37.7\,\% in the long-horizon (type 2) families; with a longer horizon a van has time for more customers after the stranded ones.
  \item \textbf{Stranded EVs} (3666 of 24250) take a van whatever their income, so they decide which vans leave the depot; the van time left over goes to free-choice customers, mostly the high-income ones.
\end{itemize}
